\documentclass[11pt]{article}
\usepackage[a4paper,margin=2.4cm]{geometry}
\usepackage{amsmath,amssymb,booktabs}
\usepackage[hidelinks]{hyperref}
\usepackage{microtype}
\setlength{\parskip}{4pt}
\title{\vspace{-1.2cm}E-HANK: what the welfare numbers are, and what they are not\vspace{-0.4cm}}
\author{Working note for Boris and Francesca}
\date{August 2026}
\begin{document}
\maketitle
\vspace{-0.6cm}

\noindent Boris asked whether the CEV numbers in the paper are subject to the well-known bias of first-order welfare evaluation (Kim and Kim, 2003; Benigno and Woodford, 2005). Short answer: the numbers \emph{are} first-order approximations to welfare, so the concern is legitimate. They are not the stochastic case in which first-order evaluation is known to give spurious rankings, but the shock is large and the adoption margin is convex, so the size of the neglected second-order terms is an empirical question, which Section~4 settles with the nonlinear solver: levels are overstated by up to 10--15\% at the headline shock, every policy ordering in the paper survives, and one sign claim (the crisis-only greening dividend) must be dropped. This note states precisely what the code computes, why it is the leading-order object, where it may fail, and what the check shows.

\section*{1. What \texttt{welfare.py} computes}

Utilitarian welfare of discount-factor type $i$ along a perfect-foresight transition indexed by shock size $\varepsilon$ is
\begin{equation}
W_i(\varepsilon)=\sum_{t\ge 0}\beta_i^{t}\Big[\int u\big(c_t(x;\varepsilon)\big)\,d\Lambda_{i,t}(x;\varepsilon)-v\big(n_t(\varepsilon)\big)\Big],
\qquad v(n)=\varphi\,\frac{n^{1+1/\nu}}{1+1/\nu},
\label{eq:W}
\end{equation}
where $x=(d,e,a)$ is the household state, $\Lambda_{i,t}$ the type-$i$ distribution, and $v(n)$ the labour disutility implied by the union's wage-setting condition (uniform hours, so it is common to all households).\footnote{With \texttt{ghh\_prefs}=0 preferences are separable; $v(n)$ never enters the household's own problem and is added ex post from the aggregate $n$ path.}

The code evaluates the transition part of \eqref{eq:W} from the \emph{linear} impulse response. \texttt{UTIL\_i} is a hetoutput ($u(c)$ aggregated over the distribution); its linear IRF, returned by \texttt{solve\_impulse\_linear}, is by construction the first derivative of aggregate flow utility with respect to the shock,
\begin{equation}
d\,\mathrm{UTIL}_{i,t}=\int u'(\bar c)\,\partial_\varepsilon c_t\,d\bar\Lambda_i+\int u(\bar c)\,\partial_\varepsilon\Lambda_{i,t},
\end{equation}
so that
\begin{equation}
dW_i \;=\; \underbrace{\frac{\mathrm{UTIL}_i^{\text{pre}}-\mathrm{UTIL}_i^{\text{base}}-dv^{\text{lvl}}}{1-\beta_i}}_{\text{standing gap between steady states (exact)}}
\;+\;\underbrace{\sum_t\beta_i^{t}\big(d\,\mathrm{UTIL}_{i,t}-v'(\bar n)\,dn_t\big)}_{\text{crisis transition (first order in }\varepsilon)},
\label{eq:dW}
\end{equation}
and, with log utility, $\chi_i=\exp\!\big((1-\beta_i)\,dW_i\big)-1$, the permanent proportional consumption change that delivers $dW_i$. The mapping from $dW_i$ to $\chi_i$ is exact; $dW_i$ itself is $O(\varepsilon)$ with the $O(\varepsilon^2)$ terms dropped. The standing term (ETS steady state versus baseline steady state) comes from two nonlinear steady-state solves and is exact.

\medskip\noindent
So the honest description of every CEV in the paper is: \emph{the first-order Taylor expansion of welfare in the shock size, around a distorted steady state, along a deterministic transition}. What is dropped is
\begin{equation}
\tfrac12\varepsilon^2 W_i''=\tfrac12\varepsilon^2\sum_t\beta_i^t\Big[\int\big(u''(\bar c)(\partial_\varepsilon c_t)^2+u'(\bar c)\,\partial^2_\varepsilon c_t\big)d\bar\Lambda_i+2\!\int u'(\bar c)\,\partial_\varepsilon c_t\,\partial_\varepsilon\Lambda_{i,t}+\int u(\bar c)\,\partial^2_\varepsilon\Lambda_{i,t}-v''\,(dn_t)^2-v'\,\partial^2_\varepsilon n_t\Big].
\label{eq:second}
\end{equation}

\section*{2. Why this is nevertheless the leading-order object}

The Kim--Kim ``spurious welfare reversal'' and the Benigno--Woodford LQ apparatus address the \emph{stochastic} problem: expected welfare $\mathbb E\,W(\varepsilon)$ under recurring shocks. There $\mathbb E[\varepsilon W']=0$ by certainty equivalence, so the welfare effect of a policy is entirely second order, $\tfrac12\sigma^2\mathbb E[W'']$, and a first-order solution gets $W''$ wrong because it omits $\partial^2_\varepsilon c$ in \eqref{eq:second}. Rankings based on it can flip. Benigno and Woodford show that when the steady state is distorted, even the LQ objective requires a second-order approximation of the constraints.

Our exercise is different in one respect that matters: it is a single deterministic MIT transition. Then $\varepsilon W_i'\neq 0$ generically---the steady state is not Pareto-efficient (incomplete markets, markups, sticky wages), and the shock is a terms-of-trade income loss with first-order distributional incidence---so the first-order term is the dominant one and \eqref{eq:dW} captures it. The linear IRF is exactly the first-order perturbation solution (Boppart, Krusell and Mitman, 2018; Auclert, Bardóczy, Rognlie and Straub, 2021), so this is the standard object one obtains from a linearised HANK; it is not the object that Kim--Kim show to be spurious. For $\varepsilon\to 0$ the levels and rankings in the paper are correct to leading order.

\section*{3. Where it can fail: the shock is not small}

Two facts in the paper qualify the ``leading order'' defence.
\begin{itemize}\itemsep2pt
\item The headline price shock doubles the world price ($\varepsilon=1$). This is not a neighbourhood of the steady state.
\item Appendix~B documents that the adoption share is strongly convex in the shock: the nonlinear-to-linear ratio of the peak $D^G$ response rises from 1.16 at $\varepsilon=1/8$ to 2.04 at $\varepsilon=3/4$, and the nonlinear solver does not converge at $\varepsilon=1$. Aggregate output is essentially linear (ratio 1.00--1.01), which is why we have described the linear solution as reliable ``for aggregate and consumption-based statistics''---a claim about allocations, not about welfare.
\end{itemize}
Orders of magnitude for the terms in \eqref{eq:second}. The curvature term $\tfrac12u''(\bar c)(dc)^2$ with consumption falling about $4$--$5\%$ per quarter is roughly $2\%$ of the first-order term---benign. The dangerous terms are those involving $\partial^2_\varepsilon c$ and $\partial^2_\varepsilon\Lambda$ through the adoption margin (switching costs paid, insulation gained, selection of who switches), which are exactly the ones the linear solution understates by up to a factor two.

The consequence is a distinction between the paper's claims:
\begin{center}\small
\begin{tabular}{lrl}
\toprule
Comparison (price shock, CEV\%) & gap & status at first order\\
\midrule
transfer $-0.82$ vs cap $-1.15$ & $0.33$ pp & probably robust; to be verified\\
flat transfer $-0.32$ vs Slutsky $-0.82$ & $0.50$ pp & probably robust; to be verified\\
cap sign flip with supply elasticity & $\sim0.7$ pp & probably robust; to be verified\\
ex-ante ETS $-1.852$ vs laissez-faire $-1.871$ & $0.02$ pp & \textbf{inside the approximation error}\\
greening dividend by persistence, $+0.01\to-0.09$ pp & $\le 0.1$ pp & \textbf{sign not established at first order}\\
\bottomrule
\end{tabular}
\end{center}
The first three rest on gaps that are themselves first order and large relative to any plausible second-order correction. The last two do not: ``ex ante barely insures'' is safe as a statement of equality, but ``the greening dividend turns negative'' (Section~6.4) is a sign claim resting on differences of a few basis points and should not be made from a first-order welfare calculation. Either we show it survives the nonlinear check or we restate it as ``zero to first order''.

\section*{4. The nonlinear check}

The sequence-space toolbox solves the nonlinear perfect-foresight transition (\texttt{solve\_impulse\_nonlinear}: Newton on the aggregate unknowns, household block re-solved exactly at every iteration). Along that path \texttt{UTIL\_i} is the exact aggregate of $u(c)$ over the exact distribution, so \eqref{eq:W} is evaluated without approximation. \texttt{run\_nl\_cev.py} does this for laissez-faire, cap, Slutsky transfer, flat transfer and the ex-ante ETS at $\varepsilon\in\{0.25,0.5,0.75\}$ (the headline $\varepsilon=1$ does not converge; the cap does not converge at $0.75$).

\begin{center}\small
\begin{tabular}{llrrrrrr}
\toprule
$\varepsilon$ & Scenario & CEV\% (linear) & CEV\% (nonlinear) & NL/L & impat. & middle & patient\\
\midrule
0.25 & Laissez-faire    & $-0.471$ & $-0.463$ & 0.98 & 0.98 & 0.98 & 0.99\\
0.25 & Ex-post cap      & $-0.288$ & $-0.285$ & 0.99 & 0.99 & 0.99 & 0.99\\
0.25 & Slutsky transfer & $-0.205$ & $-0.198$ & 0.97 & 0.95 & 0.97 & 0.99\\
0.25 & Flat transfer    & $-0.079$ & $-0.076$ & 0.96 & 1.16 & 0.94 & 0.99\\
0.25 & Ex-ante ETS      & $-0.434$ & $-0.428$ & 0.99 & 0.99 & 0.99 & 0.99\\
\midrule
0.5  & Laissez-faire    & $-0.940$ & $-0.900$ & 0.96 & 0.95 & 0.96 & 0.97\\
0.5  & Ex-post cap      & $-0.575$ & $-0.565$ & 0.98 & 0.97 & 0.98 & 0.99\\
0.5  & Slutsky transfer & $-0.409$ & $-0.377$ & 0.92 & 0.88 & 0.92 & 0.96\\
0.5  & Flat transfer    & $-0.159$ & $-0.139$ & 0.87 & 1.46 & 0.83 & 0.97\\
0.5  & Ex-ante ETS      & $-0.909$ & $-0.875$ & 0.96 & 0.97 & 0.97 & 0.95\\
\midrule
0.75 & Laissez-faire    & $-1.407$ & $-1.301$ & 0.93 & 0.92 & 0.93 & 0.93\\
0.75 & Slutsky transfer & $-0.613$ & $-0.524$ & 0.86 & 0.79 & 0.85 & 0.91\\
0.75 & Flat transfer    & $-0.238$ & $-0.175$ & 0.74 & 1.89 & 0.65 & 0.92\\
0.75 & Ex-ante ETS      & $-1.382$ & $-1.283$ & 0.93 & 0.94 & 0.93 & 0.90\\
\bottomrule
\end{tabular}
\end{center}

\noindent Four things follow.
\begin{enumerate}\itemsep2pt
\item \emph{Levels.} The first-order CEV overstates the welfare loss, and the gap grows with the shock: 1--2\% at $\varepsilon=0.25$, 4--13\% at $0.5$, 7--26\% at $0.75$. The direction is the one the convexity of the adoption margin predicts: the true switching wave is larger (nonlinear/linear peak $D^G$ of 1.3, 1.7, 2.0 at the three sizes) and cushions more. Extrapolating, the paper's levels at $\varepsilon=1$ are upper bounds on the losses, by something like 10--15\% for laissez-faire and the cap and more for the transfers. They should be reported with that caveat.
\item \emph{Ordering.} Flat $>$ Slutsky $>$ cap $>$ ETS $>$ laissez-faire at every size where all scenarios converge, linear and nonlinear alike. The first-order gaps barely move: Slutsky$-$cap $+0.083\to+0.087$ pp at $0.25$ and $+0.166\to+0.188$ at $0.5$; flat$-$Slutsky $+0.125\to+0.122$ and $+0.250\to+0.238$. The paper's welfare rankings stand.
\item \emph{Ex ante vs laissez-faire.} The gap is $+0.037\to+0.034$ pp at $0.25$, $+0.031\to+0.025$ at $0.5$, $+0.025\to+0.018$ at $0.75$: a few basis points, equal to the standing gain of the ETS steady state ($+0.043$) minus a crisis-time dividend that is small and shrinking in $\varepsilon$. ``Ex ante barely insures'' is robust. The sign of the crisis-only dividend is not something a first-order calculation can establish, and Section~6.4 should be rephrased accordingly (``zero to first order, and if anything falling with the size and persistence of the shock'').
\item \emph{Distribution.} The one place the linear solution errs in the other direction is the impatient type under the flat transfer: NL/L of $1.16$, $1.46$, $1.89$. The linear solution understates the gain of high-MPC households when the transfer relaxes a binding liquidity constraint (asymmetry of $u'$). The distributional result of Section~5.5 (flat dominates energy-indexed targeting) is reinforced, not weakened, by the nonlinear check.
\end{enumerate}

On precedents: we do not know of a HANK paper solved with the sequence-space Jacobian that derives a second-order welfare objective. Bhandari, Evans, Golosov and Sargent (2021) do second-order (small-noise) welfare in HANK, outside SSJ, for optimal Ramsey policy under recurring shocks---the stochastic case. Dávila and Schaab (2025) provide the exact decomposition of heterogeneous-agent welfare changes and would be the natural frame if we go further than the ordering check.

\vspace{0.4cm}\small
\noindent\textbf{References.}
Auclert, A., B.~Bardóczy, M.~Rognlie, L.~Straub (2021), ``Using the Sequence-Space Jacobian to Solve and Estimate Heterogeneous-Agent Models,'' \emph{Econometrica} 89(5), 2375--2408.
Benigno, P., M.~Woodford (2005), ``Inflation Stabilization and Welfare: The Case of a Distorted Steady State,'' \emph{Journal of the European Economic Association} 3(6), 1185--1236.
Benigno, P., M.~Woodford (2012), ``Linear-Quadratic Approximation of Optimal Policy Problems,'' \emph{Journal of Economic Theory} 147(1), 1--42.
Bhandari, A., D.~Evans, M.~Golosov, T.~Sargent (2021), ``Inequality, Business Cycles, and Monetary-Fiscal Policy,'' \emph{Econometrica} 89(6), 2559--2599.
Boppart, T., P.~Krusell, K.~Mitman (2018), ``Exploiting MIT Shocks in Heterogeneous-Agent Economies,'' \emph{Journal of Economic Dynamics and Control} 89, 68--92.
Dávila, E., A.~Schaab (2025), ``Welfare Assessments with Heterogeneous Individuals,'' \emph{Journal of Political Economy} 133(9), 2918--2961 (NBER WP 30571, 2022).
Kim, J., S.~H.~Kim (2003), ``Spurious Welfare Reversals in International Business Cycle Models,'' \emph{Journal of International Economics} 60(2), 471--500.
Schmitt-Grohé, S., M.~Uribe (2004), ``Solving Dynamic General Equilibrium Models Using a Second-Order Approximation to the Policy Function,'' \emph{Journal of Economic Dynamics and Control} 28(4), 755--775.
\end{document}
